\chapter{2020年新高考I卷}

\section{单选题}

\begin{problem}
设集合 \(A=\{x\mid 1\le x\le 3\}\),\(B=\{x\mid 2<x<4\}\),则 \(A\cup B\) 等于(\quad)

\choices
    {\(\{x\mid 2<x\le 3\}\)}
    {\(\{x\mid 2\le x\le 3\}\)}
    {\(\{x\mid 1\le x<4\}\)}
    {\(\{x\mid 1<x<4\}\)}
\end{problem}

\begin{problem}
\(\dfrac{2-\i}{1+2\i}\) 等于(\quad)

\choices
    {\(1\)}
    {\(-1\)}
    {\(\i\)}
    {\(-\i\)}
\end{problem}

\begin{problem}
\(6\) 名同学到甲,乙,丙三个场馆做志愿者,每名同学只去 \(1\) 个场馆,甲场馆安排 \(1\) 名,乙场馆安排 \(2\) 名,丙场馆安排 \(3\) 名,则不同的安排方法共有(\quad)

\choices
    {\(120\) 种}
    {\(90\) 种}
    {\(60\) 种}
    {\(30\) 种}
\end{problem}

\begin{problem}
日晷是中国古代用来测定时间的仪器,利用与晷面垂直的晷针投射到晷面的影子来测定时间. 把地球看成一个球(球心记为 \(O\)),地球上一点 \(A\) 的纬度是指 \(OA\) 与地球赤道所在平面所成角,点 \(A\) 处的水平面是指过点 \(A\) 且与 \(OA\) 垂直的平面. 在点 \(A\) 处放置一个日晷,若晷面与赤道所在平面平行,点 \(A\) 处的纬度为北纬 \(40^\circ\),则晷针与点 \(A\) 处的水平面所成角为(\quad)

\begin{center}
\bitmapfigure[max width=0.42\linewidth]{img/2020/new_gaokao_paper_1/q4_fig1.png}
\end{center}

\choices
    {\(20^\circ\)}
    {\(40^\circ\)}
    {\(50^\circ\)}
    {\(90^\circ\)}
\end{problem}

\begin{problem}
某中学的学生积极参加体育锻炼,其中有 \(96\%\) 的学生喜欢足球或游泳,\(60\%\) 的学生喜欢足球,\(82\%\) 的学生喜欢游泳,则该中学既喜欢足球又喜欢游泳的学生数占该校学生总数的比例是(\quad)

\choices
    {62\%}
    {56\%}
    {46\%}
    {42\%}
\end{problem}

\begin{problem}
基本再生数 \(R_0\) 与世代间隔 \(T\) 是新冠肺炎的流行病学基本参数. 基本再生数指一个感染者传染的平均人数,世代间隔指相邻两代间传染所需的平均时间. 在新冠肺炎疫情初始阶段,可以用指数模型:\(I(t)=\e^{rt}\) 描述累计感染病例数 \(I(t)\) 随时间 \(t\)(单位:天)的变化规律,指数增长率 \(r\) 与 \(R_0\),\(T\) 近似满足 \(R_0=1+rT\). 有学者基于已有数据估计出 \(R_0=3.28\),\(T=6\). 据此,在新冠肺炎疫情初始阶段,累计感染病例数增加 1 倍需要的时间约为(\(\ln 2\approx 0.69\))(\quad)

\choices
    {\(1.2\) 天}
    {\(1.8\) 天}
    {\(2.5\) 天}
    {\(3.5\) 天}
\end{problem}

\begin{problem}
已知 \(P\) 是边长为 2 的正六边形 \(ABCDEF\) 内的一点,则 \(\overrightarrow{AP}\cdot\overrightarrow{AB}\) 的取值范围是(\quad)

\choices
    {\(\left( -2,6 \right)\)}
    {\(\left( -6,2 \right)\)}
    {\(\left( -2,4 \right)\)}
    {\(\left( -4,6 \right)\)}
\end{problem}

\begin{problem}
若定义在 \(\R\) 上的奇函数 \(f(x)\) 在 \((-\infty,0)\) 上单调递减,且 \(f(2)=0\),则满足 \(xf(x-1)\ge 0\) 的 \(x\) 的取值范围是(\quad)

\choices
    {\([ -1,1]\cup[3,+\infty)\)}
    {\([ -3,-1]\cup[0,1]\)}
    {\([ -1,0]\cup[1,+\infty)\)}
    {\([ -1,0]\cup[1,3]\)}
\end{problem}

\section{多选题}

\begin{problem}
已知曲线 \(C:mx^2+ny^2=1\). (\quad)

\choices
    {若 \(m>n>0\),则 \(C\) 是椭圆,其焦点在 \(y\) 轴上}
    {若 \(m=n>0\),则 \(C\) 是圆,其半径为 \(\sqrt{n}\)}
    {若 \(mn<0\),则 \(C\) 是双曲线,其渐近线方程为 \(y=\pm\sqrt{-\dfrac{m}{n}}x\)}
    {若 \(m=0\),\(n>0\),则 \(C\) 是两条直线}
\end{problem}

\begin{problem}
如图是函数 \(y=\sin(\omega x+\varphi)\) 的部分图象,则 \(\sin(\omega x+\varphi)\) 等于(\quad)

\begin{center}
\bitmapfigure[max width=0.42\linewidth]{img/2020/new_gaokao_paper_1/q10_fig1.png}
\end{center}

\choices
    {\(\sin\left( x+\dfrac{\pi}{3} \right)\)}
    {\(\sin\left( \dfrac{\pi}{3}-2x \right)\)}
    {\(\cos\left( 2x+\dfrac{\pi}{6} \right)\)}
    {\(\cos\left( \dfrac{5\pi}{6}-2x \right)\)}
\end{problem}

\begin{problem}
已知 \(a>0\),\(b>0\),且 \(a+b=1\),则(\quad)

\choices
    {\(a^2+b^2\ge\dfrac{1}{2}\)}
    {\(2^{a-b}>\dfrac{1}{2}\)}
    {\(\log_2 a+\log_2 b\ge -2\)}
    {\(\sqrt{a}+\sqrt{b}\le\sqrt{2}\)}
\end{problem}

\begin{problem}
信息熵是信息论中的一个重要概念. 设随机变量 \(X\) 所有可能的取值为 \(1\),\(2\),\(\cdots\),\(n\),且 \(P(X=i)=p_i>0\)(\(i=1\),\(2\),\(\cdots\),\(n\)),\(\sum_{i=1}^{n}p_i=1\),定义 \(X\) 的信息熵 \(H(X)=-\sum_{i=1}^{n}p_i\log_2 p_i\). (\quad)

\choices
    {若 \(n=1\),则 \(H(X)=0\)}
    {若 \(n=2\),则 \(H(X)\) 随着 \(p_1\) 的增大而增大}
    {若 \(p_i=\dfrac{1}{n}\)(\(i=1\),\(2\),\(\cdots\),\(n\)),则 \(H(X)\) 随着 \(n\) 的增大而增大}
    {若 \(n=2m\),随机变量 \(Y\) 所有可能的取值为 \(1\),\(2\),\(\cdots\),\(m\),且 \(P(Y=j)=p_j+p_{2m+1-j}\)(\(j=1\),\(2\),\(\cdots\),\(m\)),则 \(H(X)\le H(Y)\)}
\end{problem}

\section{填空题}

\begin{problem}
斜率为 \(\sqrt{3}\) 的直线过抛物线 \(C:y^2=4x\) 的焦点,且与 \(C\) 交于 \(A\),\(B\) 两点,则 \(|AB|=\fillinblank{}\).
\end{problem}

\begin{problem}
将数列 \(\{2n-1\}\) 与 \(\{3n-2\}\) 的公共项从小到大排列得到数列 \(\{a_n\}\),则 \(\{a_n\}\) 的前 \(n\) 项和为\fillinblank{}.
\end{problem}

\begin{problem}
某中学开展劳动实习,学生加工制作零件,零件的截面如图所示. \(O\) 为圆孔及轮廓圆弧 \(AB\) 所在圆的圆心,\(A\) 是圆弧 \(AB\) 与直线 \(AG\) 的切点,\(B\) 是圆弧 \(AB\) 与直线 \(BC\) 的切点,四边形 \(DEFG\) 为矩形,\(BC\perp DG\),垂足为 \(C\),\(\tan\angle ODC=\dfrac{3}{5}\),\(BH\parallel DG\),\(EF=12 \mathrm{~cm}\),\(DE=2 \mathrm{~cm}\),\(A\) 到直线 \(DE\) 和 \(EF\) 的距离均为 \(7 \mathrm{~cm}\),圆孔半径为 \(1 \mathrm{~cm}\),则图中阴影部分的面积为\fillinblank{}\ \(\mathrm{cm}^2\).

\FigureLayoutDeclare{img/2020/new_gaokao_paper_1/q15_fig1.png}{12.0cm}{24bp 49bp 14bp 14bp}
\bitmapfigure[max width=12.0cm]{img/2020/new_gaokao_paper_1/q15_fig1.png}
\end{problem}

\begin{problem}
已知直四棱柱 \(ABCD-A_1B_1C_1D_1\) 的棱长均为 2,\(\angle BAD=60^\circ\). 以 \(D_1\) 为球心,\(\sqrt{5}\) 为半径的球面与侧面 \(BCC_1B_1\) 的交线长为\fillinblank{}.
\end{problem}

\section{解答题}

\begin{problem}
在条件
\begin{circlelist}
\item \(ac=\sqrt{3}\),
\item \(c\sin A=3\),
\item \(c=\sqrt{3}b\)
\end{circlelist}
这三个条件中任选一个,补充在下面问题中,若问题中的三角形存在,求 \(c\) 的值;若问题中的三角形不存在,说明理由.

问题:是否存在 \(\triangle ABC\),它的内角 \(A\),\(B\),\(C\) 的对边分别为 \(a\),\(b\),\(c\),且 \(\sin A=\sqrt{3}\sin B\),\(C=\dfrac{\pi}{6}\),\fillinblank{}?
\end{problem}

\begin{problem}
已知公比大于 1 的等比数列 \(\{a_n\}\) 满足 \(a_2+a_4=20\),\(a_3=8\).

\begin{enumerate}
\item 求 \(\{a_n\}\) 的通项公式;
\item 记 \(b_m\) 为 \(\{a_n\}\) 在区间 \((0,m]\)(\(m\in\N^*\))中的项的个数,求数列 \(\{b_m\}\) 的前 100 项和 \(S_{100}\).
\end{enumerate}
\end{problem}

\begin{problem}
为加强环境保护,治理空气污染,环境监测部门对某市空气质量进行调研,随机抽查了 \(100\) 天空气中的 \(\mathrm{PM}_{2.5}\) 和 \(\mathrm{SO}_2\) 浓度(单位:\(\mu\mathrm{g}/\mathrm{m}^3\)),得下表:

\begin{center}
\begin{tabular}{c|ccc}
\(\mathrm{PM}_{2.5}\backslash\mathrm{SO}_2\) & \([0,50]\) & \((50,150]\) & \((150,475]\) \\
\hline
\([0,35]\) & 32 & 18 & 4 \\
\((35,75]\) & 6 & 8 & 12 \\
\((75,115]\) & 3 & 7 & 10 \\
\end{tabular}
\end{center}

\begin{enumerate}
\item 估计事件``该市一天空气中 \(\mathrm{PM}_{2.5}\) 浓度不超过 75,且 \(\mathrm{SO}_2\) 浓度不超过 150''的概率;
\item 根据所给数据,完成下面的 \(2\times 2\) 列联表:

\begin{center}
\begin{tabular}{c|cc}
\(\mathrm{PM}_{2.5}\backslash\mathrm{SO}_2\) & \([0,150]\) & \((150,475]\) \\
\hline
\([0,75]\) & & \\
\((75,115]\) & & \\
\end{tabular}
\end{center}

\item 根据(2)中的列联表,判断是否有 99\% 的把握认为该市一天空气中 \(\mathrm{PM}_{2.5}\) 浓度与 \(\mathrm{SO}_2\) 浓度有关?
\end{enumerate}

附:\(K^2=\dfrac{n(ad-bc)^2}{(a+b)(c+d)(a+c)(b+d)}\).

\begin{center}
\begin{tabular}{c|ccc}
\(P(K^2\ge k)\) & 0.050 & 0.010 & 0.001 \\
\hline
\(k\) & 3.841 & 6.635 & 10.828 \\
\end{tabular}
\end{center}
\end{problem}

\begin{problem}
如图,四棱锥 \(P-ABCD\) 的底面为正方形,\(PD\perp\)底面 \(ABCD\). 设平面 \(PAD\) 与平面 \(PBC\) 的交线为 \(l\).

\begin{enumerate}
\item 证明:\(l\perp\)平面 \(PDC\);
\item 已知 \(PD=AD=1\),\(Q\) 为 \(l\) 上的点,求 \(PB\) 与平面 \(QCD\) 所成角的正弦值的最大值.
\end{enumerate}

\bitmapfigure[float=inline,max width=0.42\linewidth]{img/2020/new_gaokao_paper_1/q20_fig1.png}
\end{problem}

\begin{problem}
已知函数 \(f(x)=a\e^{x-1}-\ln x+\ln a\).

\begin{enumerate}
\item 当 \(a=\e\) 时,求曲线 \(y=f(x)\) 在点 \(\left( 1,f(1) \right)\) 处的切线与两坐标轴围成的三角形的面积;
\item 若 \(f(x)\ge 1\),求 \(a\) 的取值范围.
\end{enumerate}
\end{problem}

\begin{problem}
已知椭圆 \(C:\dfrac{x^2}{a^2}+\dfrac{y^2}{b^2}=1\)(\(a>b>0\))的离心率为 \(\dfrac{\sqrt{2}}{2}\),且过点 \(A\left( 2,1 \right)\).

\begin{enumerate}
\item 求 \(C\) 的方程;
\item 点 \(M\),\(N\) 在 \(C\) 上,且 \(AM\perp AN\),\(AD\perp MN\),\(D\) 为垂足. 证明:存在定点 \(Q\),使得 \(|DQ|\) 为定值.
\end{enumerate}
\end{problem}

